\subsection{The exact identity-relative CM family}

\begin{result}{E.201}{Proved}
{A full gentle element with identity relative action and the original unequal CM frames}
Use the literal modular marking
\[
 \mathfrak B=C_2\amalg C_3=\langle s,r\mid s^2=r^3=1\rangle,
 \quad \tau=sr,\quad
 F=\ker(\mathfrak B\to S_3),\quad
 x=\tau^{-2},\quad y=sxs,
\]
and put \(H=\overline{F'F^3}\), \(R_H=R_{\rm nil}(H)\).
There exists a full positive gentle pair \(g=(1,f)\) whose
normalized modular automorphism \(\Phi\) satisfies
\[
 \begin{gathered}
 f\in R_H\subset H'\subset F',\qquad
 \Phi\equiv\operatorname{id}\pmod {R_H},\qquad
 \Phi(\tau)=\tau,\\
 \Phi(s)=\beta s\beta^{-1},\qquad
 \Phi(r)=ArA^{-1},\qquad
 \beta\in R_H,\quad A\in H.
 \end{gathered}
\tag{E.201.1}
\]
For the actual modulus-twelve quotient
\[
 \begin{gathered}
 q:\mathfrak B\longrightarrow
 Q=\operatorname{SL}_2(\mathbb Z/12)/\{\pm I\},\\
 q(s)=[S],\quad S=\begin{pmatrix}0&1\\-1&0\end{pmatrix},
 \qquad
 q(r)=[R],\quad R=\begin{pmatrix}0&-1\\1&-1\end{pmatrix},
 \end{gathered}
\tag{E.201.2}
\]
its unique pure accessory frames have
\[
                    q(A)=[5I],\qquad q(\beta)=1.
\tag{E.201.3}
\]
In the signed convention of E.95 and E.181, the complete
Kummer shear is \(\rho_g=0\), the congruence intertwiner is
\(V_g=I\), and
\[
                  T_g=5I,\qquad W_g=I,
                  \qquad (\epsilon_3(T_g),\epsilon_2(W_g))=(1,0).
\tag{E.201.4}
\]
In particular \(\Delta(g)=1\).  This reconstructs the exact
identity-relative family of restored LM7 E.58, including its
literal full-source frames.  It is distinct from the compatible
involution construction E.97.  No order, involution, arithmetic
membership, or value of an individual tripod test is asserted
for the full lift.
\end{result}

\begin{proof}
\emph{The relative product retains its finite base.}
The six epimorphisms \(C_2\amalg C_3\to S_3\) send \(s\)
to a transposition and \(r\) to a three-cycle.  They form one
orbit under \(\operatorname{Aut}(S_3)\), so have one kernel.
Consequently \(F\) is characteristic in \(\mathfrak B\).
The usual six-coset Schreier calculation gives its actual free
profinite basis \(x,y\); thus \(H\) is characteristic and
\[
                 D=\mathfrak B/H,\qquad |D|=6\cdot9=54.
\tag{E.201.5}
\]
For each prime \(p\), write
\(\Gamma_p=\mathfrak B/O^p(H)\), with its specified map to
\(D\).  There is a canonical isomorphism of actual source groups
\[
             \mathfrak B/R_H\simeq\prod_{p,D}\Gamma_p.
\tag{E.201.6}
\]
Indeed \(H/R_H\simeq\prod_pH^{(p)}\): finite nilpotent
quotients are products of their Sylow subgroups, and the map
onto every finite product of distinct prime-power quotients
is onto because they have no nontrivial common quotient.
Compactness gives the full product.  To lift a tuple in the
right side of E.201.6, first lift its common \(D\)-value to
\(\mathfrak B\), then correct the remaining tuple in
\(\prod_pH^{(p)}\).  This proves the claimed source comparison
without choosing a splitting over \(D\).

\emph{The finite matrix and Fermat images.}
Let
\[
 V_4=\left\{
           \begin{pmatrix}1&2a\\2b&1\end{pmatrix}
                  :a,b\in\mathbb F_2\right\}
          \subset\operatorname{SL}_2(\mathbb Z/4).
\tag{E.201.7}
\]
It is a four-group not containing \(-I\).  The literal images are
\[
 q(x)=\left[\begin{pmatrix}1&2\\0&1\end{pmatrix}\right],
 \qquad
 q(y)=\left[\begin{pmatrix}1&0\\-2&1\end{pmatrix}\right].
\tag{E.201.8}
\]
Their subgroup is the entire kernel of reduction \(Q\to S_3\).
It has order \(96\), and each of its elements has a unique
signed representative whose reduction at four belongs to \(V_4\).
The Chinese remainder coordinates give the actual group isomorphism
\[
               q(F)\simeq V_4\times\operatorname{SL}_2(\mathbb F_3).
\tag{E.201.9}
\]
For an explicit verification, enumeration from the two matrices in
E.201.2 gives all \(576\) projective determinant-one matrices;
enumeration from E.201.8 gives all \(96\) elements of the
displayed product.  All \(96^2\) products of their unique signed
representatives obey E.201.9.  The standalone exact certificate
below performs these finite enumerations from the literal generators.

In \(\operatorname{SL}_2(\mathbb F_3)\), the matrices
\[
 i=\begin{pmatrix}0&1\\2&0\end{pmatrix},\qquad
 j=\begin{pmatrix}1&1\\1&2\end{pmatrix}
\]
satisfy \(i^2=j^2=-I\), \(ij=-ji\), and generate \(Q_8\).
This normal subgroup has quotient of order three.  The literal
commutators generate \(Q_8\), whereas cubes retain \(V_4\)
and kill the three-quotient.  Therefore
\[
                  q(H)=V_4\times Q_8,
                     \qquad |q(H)|=32.
\tag{E.201.10}
\]
In particular \(q\) kills \(O^2(H)\).

Retain both actual finite coordinates and set
\[
 J=\operatorname{im}(\mathfrak B\to Q\times D)
     =Q\times_{Q/q(H)}D.
\tag{E.201.11}
\]
To prove the equality, lift a proposed \(D\)-value; the
discrepancy in \(Q\) lies in \(q(H)\), and an element of
\(H\) corrects it.  Thus \(|J|=32\cdot54=1728\).
It is a quotient of \(\Gamma_2\), and the kernel of
\(\Gamma_2\to J\) is pro-two because it lies in \(H^{(2)}\).
An independent affine model of \(D\), with multiplication
\((L,v)(M,w)=(LM,v+Lw)\), is given by
\[
 s\mapsto\left(\begin{pmatrix}0&1\\1&0\end{pmatrix},\binom00\right),
 \qquad
 r\mapsto\left(\begin{pmatrix}0&2\\1&2\end{pmatrix},\binom01\right).
\tag{E.201.12}
\]
The torsion relations hold, and \(x,y\) become the two standard
translations.  The generated image has order \(54\); the
kernel is consequently exactly \(H\).  The certificate constructs
the \(1728\)-element joint group from these same generator pairs,
so its two finite coordinates are not identified abstractly after
the computation.

\emph{Lifting a prescribed centralizer coordinate.}
With the convention \([x,y]=xyx^{-1}y^{-1}\), direct multiplication
gives
\[
 q([x,y])=\left[\begin{pmatrix}1&8\\8&5\end{pmatrix}\right],
 \qquad q([x,y]^2)=[5I]=:z_0.
\tag{E.201.13}
\]
The element \(z_0\) is central of order two in \(Q\); its
coordinates in E.201.9 are \((I,-I)\).  Since the word
\([x,y]^2\) belongs to \(H\), its joint value is
\((z_0,1)\in J\).  This value centralizes the actual image of
\(r\) in \(J\).

We need a centralizer in the full relative pro-two source, not
only in this finite group.  If \(G\to E\) has pro-two kernel
\(K\) and \(a\in G\) has order three with image of order three,
then
\[
                    C_G(a)\longrightarrow C_E(\bar a)
                    \quad\hbox{is onto}.
\tag{E.201.14}
\]
Given a lift \(g\) of a specified element of the target
centralizer, \(\langle a\rangle\) and
\(\langle gag^{-1}\rangle\) are complements to \(K\) in
the same preimage of \(\langle\bar a\rangle\).  Finite
coprime complement conjugacy adjusts them by an element of
the kernel.  In this special case that finite assertion can
also be proved by induction on the two-kernel, using a
nonzero invariant elementary abelian subgroup of its centre
and averaging over the group of order three to kill its
first cohomology.  The closed sets of such conjugators in
\(K\) have the finite intersection property; compactness
gives the adjustment in the full pro-two kernel.  Because
the two chosen generators have the same nontrivial image,
the adjustment conjugates the specified generator itself.
This proves E.201.14.

Apply it to \(\Gamma_2\to J\) and \(a=r\).  The retained
\(S_3\) coordinate detects the order of \(r\).  We obtain
\[
 c\in C_{\Gamma_2}(r)\cap H^{(2)},
                 \qquad c\mapsto(z_0,1)\text{ in }J.
\tag{E.201.15}
\]
Thus \(c\) can be chosen of the form \(u[x,y]^2\), where
\(u\) lies in the pro-two kernel of the actual joint map.

\emph{The full peripheral lift.}
Use the identity automorphism in every factor of E.201.6,
with specified peripheral conjugators
\[
 b_s=1,\qquad b_\tau=1,\qquad
 b_r=(c\text{ in the two-primary factor};\ 1\text{ elsewhere}).
\tag{E.201.16}
\]
They have one common value in \(D\).  The last tuple centralizes
\(r\) in each component, so these are legitimate literal
conjugators for the identity automorphism of
\(\mathfrak B/R_H\).  Apply the already independently proved
ordinary lifting theorem E.57 with this exact subgroup \(H\).
Its hypothesis \(\overline{[R_H,\mathfrak B]}=R_H\) is part
of that proof.  It yields a full peripheral automorphism
retaining all three tuples in E.201.16.

Its \(\tau\)-conjugator belongs to \(R_H\).  Conjugating
the automorphism by its inverse fixes \(\tau\) literally and
preserves every relative coordinate.  Call the resulting
automorphism \(\Phi\), with conjugators \(\beta,A\) at
\(s,r\).  Then \(\beta\in R_H\), \(A\in H\), and their
finite values are precisely E.201.3.  Also \(\Phi\) is the
identity modulo \(R_H\).  Therefore \(f=s\Phi(s)\) belongs
to \(R_H\), which lies in \(H'\) because the abelian quotient
of \(H\) is already pronilpotent.

The actual torsion and product relations now give
\[
 \begin{gathered}
 \Phi(s)=sf,\quad \Phi(r)=f^{-1}r,\quad
 \Phi(x)=x,\quad \Phi(y)=f^{-1}yf,\\
 f(y,x)f(x,y)=1,\qquad
 f((xy)^{-1},x)f(y,(xy)^{-1})f(x,y)=1.
 \end{gathered}
\tag{E.201.17}
\]
For the last identity use the literal action of \(r\), which
cycles \(x,y,(xy)^{-1}\).  Since \(F\) is characteristic,
the restriction of \(\Phi\) is an actual automorphism of
the full free profinite group.  Together with \(f\in F'\),
E.201.17 proves full membership \((1,f)\in\widehat{\mathrm{GT}}_0\).

\emph{The unique full frames and the signed conclusion.}
The free profinite product centralizer theorem gives
\[
 C_{\mathfrak B}(s)=\langle s\rangle,\qquad
 C_{\mathfrak B}(r)=\langle r\rangle
\]
\citep[Theorem B', \S4, pp.159--160]{SymHerfortRibes1985}.
Each intersects \(F\) trivially, as seen in the specified
\(S_3\) quotient.  Hence the conjugators \(\beta,A\in F\)
are the unique pure frames of this full normalized automorphism.
The construction has not asserted that the nontrivial relative
centralizer element lifts to a full centralizer element; it lifts
as the frame of the new full automorphism.

Because \(\beta\in R_H\subset F'\), E.181's exact identity
\((\beta^{-1})^{\rm ab}=(\rho_g,\rho_g)\) gives
\(\rho_g=0\) in \(\widehat{\mathbb Z}\), not just modulo
three.  Thus \(V_g=U^{8\rho_g}=I\).  The signed pure
representation at twelve has the unique representatives used
in E.201.9.  The lift of \([5I]\) with mod-four component in
\(V_4\) is exactly \(5I\); the lift of the identity is \(I\).
Taking inverse frames therefore gives \(T_g=5I\), \(W_g=I\).
For the \emph{positive} matrix \(R\) in E.201.2, the line
\(\ker(R-I)\) modulo three is generated by \((1,2)^t\).
The first frame sends this vector to \((2,1)^t\), its negative;
the second frame is the identity.  This proves E.201.4 and,
by the marked CM comparison E.95, \(\Delta(g)=1\).
In particular the full lift is nontrivial.  By E.191 it is
outside the full GT subgroup as well as the arithmetic subgroup.
The ordinary lifting theorem imposes no additional equivariance
or order equation, and this conclusion does not supply one.
\end{proof}

The independently written certificate
\path{code/ghost_audit_relative_cm.py} and its saved
\path{code/ghost_audit_relative_cm.out} and JSON output verify the
orders \(576,96,32,54,1728\), all \(9216\) products in the
signed decomposition, the exact joint fibres, the literal
commutator square, and both signed CM values.  Projective
canonical representatives are deliberately separated from the
positive \(R\)-matrix when testing its fixed line.  The infinite
centralizer and peripheral lifts are the proofs above; finite
enumeration is not substituted for them.
